\documentclass[11pt,a4paper]{article}

\usepackage[utf8]{inputenc}
\usepackage[T1]{fontenc}
\usepackage{amsmath,amssymb,amsfonts,amsthm}
\usepackage{booktabs}
\usepackage{array}
\usepackage{tabularx}
\usepackage{geometry}
\usepackage{hyperref}
\usepackage{xcolor}
\usepackage{enumitem}
\geometry{margin=25mm}
\hypersetup{colorlinks=true,linkcolor=blue!60!black,urlcolor=blue!60!black,citecolor=blue!60!black}

\newtheorem{proposition}{Proposition}
\newtheorem{remark}{Remark}
\newcommand{\R}{\mathbb{R}}
\newcommand{\bq}{\mathbf{q}}
\newcommand{\bv}{\mathbf{v}}
\newcommand{\ba}{\mathbf{a}}
\newcommand{\bu}{\mathbf{u}}
\newcommand{\bx}{\mathbf{x}}
\newcommand{\by}{\mathbf{y}}
\newcommand{\bh}{\mathbf{h}}
\newcommand{\bM}{\mathbf{M}}
\newcommand{\bY}{\mathbf{Y}}
\newcommand{\bB}{\mathbf{B}}
\newcommand{\bJ}{\mathbf{J}}
\newcommand{\bA}{\mathbf{A}}
\newcommand{\bK}{\mathbf{K}}
\newcommand{\bP}{\mathbf{P}}
\newcommand{\bS}{\mathbf{S}}
\newcommand{\bW}{\mathbf{W}}
\newcommand{\bpi}{\boldsymbol{\pi}}
\newcommand{\btheta}{\boldsymbol{\theta}}
\newcommand{\bxi}{\boldsymbol{\xi}}
\newcommand{\blambda}{\boldsymbol{\lambda}}
\newcommand{\norm}[1]{\left\lVert #1 \right\rVert}

\title{MOSAIC: Model-Aware Simultaneous Anthropometry, Inertia and Control\\
Estimation From Repeated Motion Captures\\[4pt]
\large Methods Reference for Forward-Dynamics Matching in UpstreamDrift}
\author{UpstreamDrift Estimation Working Reference}
\date{Revision 2026-10-05 \quad (separate research and experiment record; see \S\ref{sec:governance})}

\begin{document}
\maketitle

\begin{abstract}
Matching a forward-dynamics model to measured kinematics requires three things
to be determined at once: the subject's segment geometry and inertial
parameters, the joint inputs that drove each trial, and a kinematic trajectory
consistent with both the observations and the equations of motion.  This
reference specifies MOSAIC, an estimator that exploits two exact structural
facts of rigid-body dynamics---the control-affine form
$\dot\bx=f(\bx)+G(\bx)\bu$ and the linearity of inverse dynamics in the
inertial parameters, $\bY(\bq,\bv,\ba)\bpi=\bB\bu+\bJ_c^{\!\top}\blambda$---to
turn the joint problem into a variable-projection Gauss--Newton method over a
physically consistent inertial manifold.  Inputs, cross-trial torque templates
and contact wrenches are eliminated exactly by a sparse linear inner solve;
kinematics, geometry and inertia are updated by a damped reduced Gauss--Newton
step whose dynamics Jacobians are exact in the inertial direction.  The method
clarifies what kinematics can and cannot identify (base-parameter
combinations, a global mass gauge, the massless-body degeneracy of
minimum-effort fitting), states where the missing information comes from in
golf (the club as a known-inertia force sensor, the unactuated base rows,
ground reactions), and places the zero-torque and zero-velocity counterfactuals
(ZTCF/ZVCF) in the same algebra as exact row and column projections of the
regressor.  A NumPy reference implementation on an analytic planar chain
recovers geometry to $0.1\,\mathrm{mm}$, observable inertial parameters to
within 4\,\% (prior 5\,\%), torques to 20\,\% RMS, and passes an open-loop
forward-replay gate (0.03\,rad RMS over 0.75\,s) in 37 outer iterations and
about five seconds for two trials, every number read from a committed,
regenerable evidence receipt.  No human capture has been qualified with the method yet; the
qualification gates and the engine integration path are specified so that
claim can be earned, not asserted.
\end{abstract}

\tableofcontents

\section{Status, Scope and Governance}
\label{sec:governance}

This document is an editable standalone research and experiment record in the
sense of \texttt{AGENTS.md} (``Modeling Reference Documentation''); it is not
the engineering design manual, whose only editable source is
\texttt{manuals/upstreamdrift}.  It is the calculation-level reference for the
MOSAIC estimator implemented in
\texttt{src/shared/python/estimation/mosaic/} and for the review of the
Dynamics-Informed Motion Estimation epic (DIME, \#11421) recorded in
\S\ref{sec:dime_review}.  It complements, and is cross-linked from, the GS3DX
record \texttt{docs/research/simscape\_matching\_reference/}.

\paragraph{Qualification Status (Truthful).}
\begin{itemize}[nosep]
  \item Verified: all equations of \S\ref{sec:structure}--\S\ref{sec:algorithm}
    against closed-form double-pendulum dynamics and synthetic ground truth
    (36 behavioural tests, \texttt{tests/unit/estimation/mosaic/}).
  \item Measured: the planar results of \S\ref{sec:results} on the analytic
    fixture (capture-A/capture-O style noise levels, no engine).
  \item Not yet done: any fit of a human capture (capture-A, capture-O), any
    3-D floating-base model, any contact-wrench estimation, any engine-backed
    regressor.  These are the child issues of the follow-up epic
    (\S\ref{sec:roadmap}).  Nothing here qualifies the GS3DX IK videos.
\end{itemize}

\section{Problem Statement and Notation}
\label{sec:problem}

A subject performs $K$ trials (swings).  Trial $k$ is observed at times
$t^k_1<\dots<t^k_{T_k}$ through markers or 2-D keypoints $\by^k_t$ with
covariance $\mathbf{R}^k_t$, optionally with ground-reaction wrenches.  A
rigid multibody model with configuration $\bq\in\mathcal{Q}$ (a Lie group when
a floating base or spherical joints are present), velocity $\bv\in\R^{n_v}$ and
acceleration $\ba$ obeys
\begin{equation}
  \bM(\bq;\bpi,s)\,\ba + \bh(\bq,\bv;\bpi,s) = \bB\bu + \bJ_c(\bq;s)^{\!\top}\blambda ,
  \label{eq:eom}
\end{equation}
with $\bh=\mathbf{C}\bv+\mathbf{g}$, actuation map $\bB\in\R^{n_v\times n_u}$
($n_u<n_v$: the floating base and any passive joints are unactuated), contact
Jacobian $\bJ_c$ and wrenches $\blambda$.  Unknowns:
\begin{center}
\begin{tabularx}{\linewidth}{@{}l l X@{}}
\toprule
Symbol & Shared/per-trial & Meaning \\
\midrule
$s\in\R^{n_s}$ & shared & geometry: segment lengths, marker attachment offsets, joint-axis offsets \\
$\bpi\in\R^{10 n_b}$ & shared & inertial parameters $(m_i,\bh_i=m_i\mathbf{c}_i,\mathbf{I}_i)$ about each body-frame origin \\
$\btheta$ & shared & unconstrained coordinates of $\bpi$ on the consistency manifold (\S\ref{sec:consistency}) \\
$\mathbf{C}^k\in\R^{C\times n_v}$ & per trial & spline coefficients of the kinematic trajectory $\bq^k(t)=\bP(t)\mathbf{C}^k$ \\
$\bu^k_t$ & per trial & net joint inputs (torques; activations via a declared actuator model) \\
$\bar\bu(\phi)$ & shared & phase-indexed torque template coupling trials \\
$\blambda^k_t$ & per trial & contact wrenches (friction-cone constrained) \\
\bottomrule
\end{tabularx}
\end{center}
$\bxi=(\mathbf{C}^1,\dots,\mathbf{C}^K,s,\btheta)$ are the \emph{outer}
variables, $w=(\bu,\bar\bu,\blambda)$ the \emph{inner} variables.  Units are SI
throughout; residuals are whitened before any norm is taken so that radians
and metres are never mixed (cf.\ ch.~21 of the AffineDrift companion).

\section{Structure Exploited}
\label{sec:structure}

\subsection{Control-Affine Form and the Inertial Regressor}
Equation~\eqref{eq:eom} is control-affine: $\dot\bx=f(\bx)+G(\bx)\bu$ with
$f=(\bv,-\bM^{-1}\bh+\bM^{-1}\bJ_c^{\!\top}\blambda)$ and
$G=(0,\bM^{-1}\bB)$.  Independently, the left-hand side of \eqref{eq:eom} is
\emph{linear} in $\bpi$ (Atkeson--An--Hollerbach; Gautier--Khalil):
\begin{equation}
  \bM(\bq)\ba+\bh(\bq,\bv) \;=\; \bY(\bq,\bv,\ba;s)\,\bpi ,
  \qquad \bY\in\R^{n_v\times 10 n_b}.
  \label{eq:regressor}
\end{equation}
Both facts hold for every rigid multibody engine (Pinocchio exposes $\bY$ as
\texttt{computeJointTorqueRegressor}; MuJoCo/MJX and Drake give it by
differentiating inverse dynamics with respect to the body inertias).  MOSAIC
consumes only $\bY$, $\bB$, $\bJ_c$ and marker forward kinematics, so it is
engine-agnostic by construction.

\subsection{ZTCF and ZVCF as Exact Projections of the Regressor}
\label{sec:ztcf}
Because $\bY$ is affine in $\ba$ and $\bh=\bY(\bq,\bv,0)\bpi$, the drift and
its parts are algebraic objects of the same regressor:
\begin{align}
  \ba_{\mathrm{ZTCF}}(t) &= -\bM(\bq)^{-1}\,\bY(\bq,\bv,0)\,\bpi
    \;+\;\bM^{-1}\bJ_c^{\!\top}\blambda_0 , \label{eq:ztcf}\\
  \ba_{\mathrm{ZVCF}}(t) &= -\bM(\bq)^{-1}\,\bY(\bq,0,0)\,\bpi
    \;+\;\bM^{-1}\bJ_c^{\!\top}\blambda_{00}, \label{eq:zvcf}\\
  \ba_{\mathrm{ZTCF}}-\ba_{\mathrm{ZVCF}} &= -\bM^{-1}\,[\bY(\bq,\bv,0)-\bY(\bq,0,0)]\,\bpi
    \quad\text{(velocity-dependent drift)} .
\end{align}
Here $\blambda_0$ ($\blambda_{00}$) are the contact reactions recomputed under
zero input (zero input and zero velocity) from the same constraint equations;
in free flight they vanish.  The measured acceleration decomposes exactly as
$\ba_{\mathrm{meas}}=\ba_{\mathrm{ZTCF}}+\bM^{-1}\bB\bu+\bM^{-1}\bJ_c^{\!\top}(\blambda-\blambda_0)$,
which is the AffineDrift pointwise decomposition $\dot\bx-f(\bx)=G(\bx)\bu$
with contact made explicit.

\paragraph{Canonical nomenclature (fixing the inconsistencies found in the
repository, \S\ref{sec:dime_review}).}
\begin{itemize}[nosep]
  \item \textbf{ZTCF}: $\bu=0$, full state kept, constraints kept, reactions
    recomputed.  \emph{Pointwise} unless the word ``trajectory'' is attached;
    the integrated object $\bx(t_0)+\int f(\bx_{\mathrm{ZTCF}}(\tau))d\tau$ is a
    different curve (it follows its own flow, not the measured state) and must
    be called the ZTCF \emph{trajectory}.
  \item \textbf{ZVCF}: $\bv=0$ \emph{and} $\bu=0$ at each instant
    (configuration-only drift: gravity plus constraint reactions).  The variant
    that keeps the control is \emph{zero-velocity, control-preserved
    acceleration} and is never abbreviated ZVCF.
  \item Passive impedance declared in the model (joint damping, ligament
    stiffness) is part of $f$; a feedback servo is part of $\bu$ and is removed
    by the ZTCF.
\end{itemize}

\subsection{Where the Information Lives: Unactuated Rows}
\label{sec:unactuated}
Let $\bS_\perp$ select the unactuated coordinates ($\bS_\perp\bB=0$).  Then
\eqref{eq:eom} gives, with \emph{no unknown input},
\begin{equation}
  \bS_\perp\,\bY(\bq,\bv,\ba)\,\bpi \;=\; \bS_\perp\bJ_c^{\!\top}\blambda .
  \label{eq:unactuated}
\end{equation}
These six-per-node (floating base) or one-per-node (passive pivot) equations
are the entire torque-free information kinematics carries about $\bpi$
(Ayusawa--Venture--Nakamura).  Equivalently, in ZTCF language,
$\bS_\perp\bM(\ba_{\mathrm{meas}}-\ba_{\mathrm{ZTCF}})=\bS_\perp\bJ_c^{\!\top}(\blambda-\blambda_0)$:
the ``torque-independent feasibility criterion'' of DIME-16 is the unactuated
projection of the regressor equation, not an additional likelihood.
The actuated rows determine $\bu$ once $\bpi$ is known; they carry information
about $\bpi$ only through priors on $\bu$.

\subsection{The Club as a Known-Inertia Force Sensor}
In golf the distal body is a manufactured object whose inertial parameters are
measurable to better than 1\,\% off-line.  With $\bpi_{\mathrm{club}}$ known
and its kinematics observed, the club's Newton--Euler equations return the
hand wrench exactly, and the unactuated rows \eqref{eq:unactuated} become
\emph{inhomogeneous} in the unknown human parameters,
$\bS_\perp\bY_{\mathrm{human}}\bpi_{\mathrm{human}}=-\bS_\perp\bY_{\mathrm{club}}\bpi_{\mathrm{club}}+\dots$,
which removes the global scale gauge without a force plate.  This is the
golf-specific anchor used in \S\ref{sec:results}; for a humanoid the analogue
is a known payload or the measured ground wrench.

\section{Identifiability}
\label{sec:identifiability}

\begin{proposition}[Scale gauge]
If no external wrench is measured and no anchor is imposed, the map
$(\bpi,\bu,\blambda)\mapsto(\alpha\bpi,\alpha\bu,\alpha\blambda)$ leaves every
residual of \eqref{eq:eom} invariant for all $\alpha>0$.  Kinematics alone
cannot identify absolute mass.
\end{proposition}

\begin{proposition}[Base parameters]
Only the row space of the stacked matrix
$[\bS_\perp\bY(\bq_t,\bv_t,\ba_t)]_t$ is observable from the unactuated rows.
A distal link's mass lumps into the proximal link's first moment and inertia
($h_{0x}+m_1 l_0$, $I_{0}+m_1 l_0^2$ in the planar chain); the orthogonal
complement is determined by the prior only.  The reference implementation
reports this subspace (\texttt{structural\_observability}) and the Fisher
information after eliminating inputs (\texttt{data\_information}).
\end{proposition}

\begin{proposition}[Massless-body degeneracy of minimum-effort fitting]
\label{prop:massless}
With free $\bpi$ constrained only by positivity, a free kinematic trajectory
and any prior that penalises input magnitude, the configuration
``all mass at the unactuated pivot, every other body massless, $\bu\equiv0$''
attains zero dynamics residual and zero effort for \emph{every} trajectory.
Without an anchor or a credible anthropometric prior the estimator converges to
it.
\end{proposition}
\begin{proof}[Sketch]
A body whose mass sits at the origin of the unactuated joint contributes
$\bS_\perp\bY\bpi=0$ for all motions (zero lever arm, zero first moment, zero
inertia); massless bodies contribute nothing; with $\bu=0$ the actuated rows
vanish identically.  Both the dynamics and the effort terms are then at their
global minimum.
\end{proof}
This degeneracy was reproduced numerically during development (\S\ref{sec:results}):
with a weak absolute prior the fit returned $m_0=2.3\,\mathrm{kg}$ (the total)
and zeros elsewhere.  Two remedies are structural, not cosmetic: the
physically consistent manifold (\S\ref{sec:consistency}), which keeps every
iterate realisable, and anchors supplying inhomogeneous information (known
club, total mass, ground wrench).  A third is honest reporting: every fit
publishes which directions of $\bpi$ came from the data.

\begin{remark}[Effort priors bias inertia downward]
Any penalty on $\norm{\bu}$ or $\norm{\Delta\bu}$ is informative about
$\bpi$ because $\bu=\bB^{+}(\bY\bpi-\bJ_c^{\!\top}\blambda)$ on the actuated
rows.  The bias is small when the unactuated rows and anchors dominate the
information (the planar result shows $\le3\,\%$) but it must be reported, and
model-free kinematic regularisation (a jerk prior on $\bq$) is preferable to
strong effort priors when acceleration noise is the problem.
\end{remark}

\section{Physically Consistent Inertia}
\label{sec:consistency}
Physical consistency is the convex cone of parameters whose $4\times4$
pseudo-inertia
$\mathcal{J}(\bpi)=\begin{bmatrix}\tfrac12\operatorname{tr}(\mathbf{I})\mathbf{1}-\mathbf{I} & \bh\\ \bh^{\!\top} & m\end{bmatrix}\succeq0$
(Wensing--Kim--Slotine).  Two representations are used:
(i) the linear 10-vector with Euclidean projection onto the cone in the
pseudo-inertia Frobenius metric (eigenvalue clipping; idempotent; identity on
consistent bodies), used to sanitise priors and report violations; and
(ii) the unconstrained log-Cholesky coordinates $\btheta$ with
$\mathcal{J}=\mathbf{L}\mathbf{L}^{\!\top}$, $\operatorname{diag}\mathbf{L}=e^{\theta}$
(Rucker--Wensing), used as outer variables so every iterate is realisable.
For planar bodies $(m,h_x,h_y,I_o)$ the cone is the rotated second-order cone
$mI_o\ge|\bh|^2$, with coordinates $(\log m,\bh,\log I_c)$,
$I_c=I_o-|\bh|^2/m$, and a closed-form projection.  Because
$\partial\bpi/\partial\btheta$ is available analytically, the dynamics
Jacobian in the inertial direction is exact: $\bW\bY\,\partial\bpi/\partial\btheta$.

\section{Estimation Problem}
\label{sec:estimation}
With trajectory basis $\bq^k_t=\bP^k_t\mathbf{C}^k$, $\bv^k_t=\mathbf{V}^k_t\mathbf{C}^k$,
$\ba^k_t=\bA^k_t\mathbf{C}^k$ (clamped quintic B-splines; position, velocity
and acceleration are all linear in $\mathbf{C}^k$), the reduced objective is
\begin{equation}
\begin{aligned}
F(\bxi)=\;&\sum_{k,t}\rho\!\left(\norm{\mathbf{R}^{-1/2}_{kt}\big[\mathbf{m}(\bq^k_t;s)-\by^k_t\big]}^2\right)
 + \min_{w}\Bigg[\sum_{k,t}\norm{\bW_\tau\big[\bY(\bq^k_t,\bv^k_t,\ba^k_t;s)\,\bpi(\btheta)-\bB\bu^k_t-\bJ_c^{\!\top}\blambda^k_t\big]}^2 \\
 &\quad + w_u^2\sum\norm{\bu^k_t}^2 + w_{\Delta u}^2\sum\norm{\Delta^2\bu^k_t}^2
   + w_T^2\sum\norm{\bu^k_t-\bar\bu_{\phi(k,t)}}^2 + w_\lambda^2\sum\norm{\blambda^k_t}^2\Bigg]\\
 &+ \norm{\bW_s(s-s_0)}^2 + \norm{\bW_\pi(\bpi(\btheta)-\bpi_0)}^2
   + \sum_j w_j^2\big(\mathbf{a}_j^{\!\top}\bpi(\btheta)-b_j\big)^2 ,
\end{aligned}
\label{eq:objective}
\end{equation}
where $\rho$ is a robust kernel (Huber/Cauchy by IRLS), $\bW_\tau=\operatorname{diag}(\gamma/\sigma_\tau)$
carries the continuation factor $\gamma\in[0,1]$ and the declared dynamics
discrepancy $\sigma_\tau$ (which must include the basis discretisation error,
\S\ref{sec:gates}), $\mathbf{a}_j^{\!\top}\bpi=b_j$ are linear anchors (total
mass, known club), and $\bar\bu$ is the shared phase-indexed template.
Everything inside $\min_w$ is linear in $w$ for fixed $\bxi$: the inner problem
is a sparse linear least-squares problem (a convex QP/SOCP when input bounds and
friction cones are active).  This is the \emph{inverse-dynamics collocation}
formulation (Remy--Thelen; Falisse \emph{et al.}): no forward integration
occurs inside the optimisation, so there is no exponential sensitivity, every
node is independent (embarrassingly parallel), and the measured states are
feasible initial iterates.

\section{Algorithm}
\label{sec:algorithm}

\subsection{Variable Projection: The ``Locally Learned Inputs''}
For fixed $\bxi$ the inner minimiser $w^\star(\bxi)$ is the solution of sparse
normal equations $(\bA^{\!\top}\bA)w=\bA^{\!\top}b$ (block-arrow structure:
per-node $\bu$ blocks, a banded $\Delta^2$ coupling, a small template block).
The inputs are therefore \emph{determined linearly from the kinematics}: this
is the exact answer to ``is the locally learned set of inputs a
linearisation?''---the inputs are a linear function of the inner problem
data, not an approximation; the only linearisation in MOSAIC is the
Gauss--Newton model of the outer step.  Eliminating $w$ exactly gives the
Golub--Pereyra/Kaufman reduced Gauss--Newton system
\begin{equation}
  \mathbf{H}_{\mathrm{red}}=\bJ_\xi^{\!\top}\big(\mathbf{I}-\bA(\bA^{\!\top}\bA)^{-1}\bA^{\!\top}\big)\bJ_\xi,
  \qquad
  \mathbf{g}_{\mathrm{red}}=\bJ_\xi^{\!\top}r^\star ,
  \label{eq:reduced}
\end{equation}
which is the Schur complement of the full Gauss--Newton matrix onto the outer
variables (bundle-adjustment structure).  The implementation verifies
\eqref{eq:reduced} against the explicit projector.

\subsection{Outer Step}
Levenberg--Marquardt on $F$ with Marquardt scaling and Nielsen damping updates
(gain ratio $\varrho$; $\mu\leftarrow\mu\max\{\tfrac13,1-(2\varrho-1)^3\}$ on
success, geometric growth on failure).  Jacobians:
observation rows analytically through $\partial\mathbf{m}/\partial\bq\cdot\bP$
and $\partial\mathbf{m}/\partial s$; dynamics rows exactly in $\btheta$
($\bW\bY\,\partial\bpi/\partial\btheta$) and by \emph{vectorised} central
differences in $(\bq,\bv,\ba,s)$---each perturbation applied to all nodes of
all trials at once, so one iteration costs $2(3n_v+n_s)$ batched regressor
evaluations plus one sparse factorisation.  Analytic derivatives (Singh--
Russell--Wensing; GRiD) replace the differences without changing the
algorithm.

\subsection{Continuation}
A schedule $\gamma\in\{0,\,10^{-2},\,10^{-1},\,1\}$ moves the problem from a
pure kinematic fit (the current IK practice, now a special case) to the fully
weighted dynamics-consistent fit.  Each stage is solved to tolerance before the
next; the measured receipts (\S\ref{sec:results}) show 8--14 iterations per
stage.  The homotopy is what makes convergence from an IK seed reliable: the
$\gamma=0$ basin is convex-like, and the dynamics terms are switched on only
once the kinematic residual is at its noise floor.

\subsection{Convergence Properties}
For zero-residual problems (synthetic truth, consistent model) Gauss--Newton
converges quadratically; with noise the rate is linear with factor governed by
residual-weighted curvature over the smallest reduced-Hessian eigenvalue.
Weakly identifiable inertial directions make that eigenvalue small; the
Levenberg--Marquardt damping with the Yamashita--Fukushima rule
$\mu\propto\norm{r}$ retains local quadratic convergence under a local error
bound even when the Jacobian is rank-deficient, which is the generic case
here.  Global convergence to a stationary point follows from the trust-region
acceptance test (Mor\'e).  Variable projection removes the inner variables'
nonlinearity from the outer iteration and is known to reduce the number of
local minima and iterations in separable problems (Golub--Pereyra 2003).

\subsection{Local Policy and Replay Acceptance}
After the fit, the discrete step of the (any-engine) model is linearised along
$(\bx^\star,\bu^\star)$ by vectorised central differences, and a backward
Riccati pass gives time-varying gains $\bK_t$.  The pair $(\bu^\star_t,\bK_t)$,
$\bu=\bu^\star_t-\bK_t(\bx-\bx^\star_t)$, is the locally optimal linear policy
around the fit---the object a whole-body controller consumes.  Acceptance is
two-sided: the \emph{open-loop} replay from one initial state is the physical
gate (a dynamically consistent fit must replay), and the closed-loop replay
reports the feedback effort needed to hold the motion, which quantifies the
fit's sensitivity.  For repeated trials the iterative-learning view
$\bu_{k+1}=\bu_k+\mathbf{L}e_k$ converges iff $\rho(\mathbf{I}-\mathbf{G}\mathbf{L})<1$;
with $\mathbf{L}$ the Riccati gain this is a cross-trial diagnostic, not the
estimator.

\subsection{Cross-Trial Torque Template}
Each trial's inputs are binned by normalised phase; the trial-balanced mean
$\bar\bu(\phi)$ enters the inner solve as a coupling prior ($w_T$) and is
reported with its principal modes (SVD of trial deviations) as the player's
motor signature.  On a planted rank-one structure the first mode explains
$>99.9\,\%$ of the variance (test).

\subsection{Complexity and Vectorisation}
Per outer iteration: $O(N)$ batched regressor evaluations ($N=\sum_k T_k$
nodes), one sparse LU of the inner normal matrix ($n_w=N n_u+Pn_u+n_\lambda$,
banded), one dense reduced system of size $n_\xi=\sum_k C_k n_v+n_s+10n_b$.
Trials enter only through the shared blocks, so the per-trial parts are
independent and batch over trials; the time axis is parallel because no
integration occurs.  The reference implementation is pure NumPy; the same
kernels accept MJX/JAX batched regressors unchanged.

\section{Qualification Gates}
\label{sec:gates}
A fit is reported with the following measured items; a missing item is
``unqualified'', never ``passed''.
\begin{enumerate}[nosep]
  \item \textbf{Kinematic whiteness.} At $\gamma=0$ the whitened marker
    residual RMS must be within a declared factor of 1 and uncorrelated in
    time.  Under-resolved bases fail this gate; dynamics weighting must not be
    switched on before it passes (Proposition~\ref{prop:massless} is reached
    otherwise, because a basis that cannot represent the acceleration makes the
    massless body the cheapest consistent explanation).
  \item \textbf{Declared discrepancy.} $\sigma_\tau$ must be at least the basis
    acceleration error times the mass scale; for the planar fixture a
    80-coefficient quintic basis over 200 nodes gives 1\,\% acceleration error.
  \item \textbf{Physical consistency.} Every reported $\bpi$ is on the cone by
    construction; violations of the sanitised prior are logged.
  \item \textbf{Observability.} Rank and per-parameter observable fraction of
    the unactuated-row regressor; posterior information after eliminating
    inputs; directions supplied by the prior are named.
  \item \textbf{Open-loop replay.} Configuration RMS and divergence time from
    one initial state with the fitted inputs; the accepted tolerance is
    declared beforehand.
  \item \textbf{Closed-loop effort.} TVLQR feedback RMS relative to the
    nominal inputs.
  \item \textbf{Cross-trial consistency.} Template explained variance and the
    spread of shared parameters across leave-one-trial-out refits.
\end{enumerate}

\section{Reference Results on the Planar Fixture}
\label{sec:results}
Fixture: pinned planar 3-link chain (pelvis--arm--club analogue), first joint
unactuated, true lengths $(0.9,0.6,0.4)\,\mathrm{m}$, two trials of 150 nodes
at 200\,Hz driven by smooth torques of RMS $0.54\,\mathrm{N\,m}$, six markers
with $0.5\,\mathrm{mm}$ noise, initial geometry $+5\,\%$, inertial prior
perturbed by $\pm10\,\%$ (projected onto the cone), anthropometric
$\sigma=15\,\%$, total mass and club inertia anchored, $\sigma_\tau=0.2\,\mathrm{N\,m}$,
$w_u=0.3$, $w_{\Delta u}=10$.  Measured values (receipt \texttt{reference\_results.json}):
\begin{center}
\begin{tabular}{@{}l r@{}}
\toprule
Initialisation (sequential IK, 2 trials) & 0.7\,s \\
Joint fit (4 stages, 8+9+9+11 iterations, all converged) & 2.3\,s \\
Replay + TVLQR (2 trials) & 1.4\,s \\
\midrule
Geometry error (identifiable lengths) & $0.024\,\mathrm{mm}$ \\
Kinematic RMS error & $0.70$, $0.74\,\mathrm{mrad}$ \\
Torque RMS error / torque RMS & $0.106/0.542 = 19.5\,\%$ \\
Inertial parameters, max relative error over observable non-zero entries & 3.6\,\% (prior: 4.6\,\%) \\
Structural observability rank & 9 of 12 ($m_0$ unobservable, as predicted) \\
Open-loop replay, configuration RMS over 0.75\,s & $0.0304$, $0.0307\,\mathrm{rad}$; no divergence \\
Closed-loop replay / feedback effort RMS & $0.00052$, $0.00053\,\mathrm{rad}$ / $0.021$, $0.030\,\mathrm{N\,m}$ \\
\bottomrule
\end{tabular}
\end{center}
Negative results retained: (a) with a 16- or 32-coefficient basis the fit
converged to the massless body (gate~1 failure); (b) with $\sigma_\tau=0.05\,\mathrm{N\,m}$
and no input smoothness the torques were dominated by acceleration noise
(RMS error $>1\,\mathrm{N\,m}$); (c) with a weak absolute prior and no club
anchor the scale gauge was resolved by the degenerate solution.  Each is now a
test or a gate.

\paragraph{Reproduction and Evidence Receipt.}
Every number in the table is read from the committed receipt
\texttt{docs/research/model\_aware\_matching/reference\_results.json}, which is
regenerated deterministically (fixed seeds) by
\begin{verbatim}
export PYTHONPATH=$PWD:$PWD/src:$PWD/src/shared/python
python3 -m src.shared.python.estimation.mosaic.reference_evidence
python3 -m pytest tests/unit/estimation/mosaic -q
\end{verbatim}
and verified by \texttt{test\_reference\_evidence.py}, which re-runs the fixture
and fails if any quoted quantity drifts from the receipt beyond its stated
tolerance (timings are recorded but never compared; they depend on the host).
Timings above were measured on NumPy~2.5/SciPy~1.18, one core of the session
container.

\section{Engine Integration Path}
\label{sec:engines}
\begin{itemize}[nosep]
  \item \textbf{Regressor provider.} Implement \texttt{RegressorModel} for
    Pinocchio (\texttt{computeJointTorqueRegressor}, batched over nodes in
    C++), MuJoCo (\texttt{mj\_inverse} differentiated w.r.t.\ body inertias,
    or MJX autodiff with \texttt{vmap}), Drake (AutoDiffXd).  Marker forward
    kinematics and Jacobians come from the same engine.
  \item \textbf{Floating base and manifold.} Replace the Euclidean spline by a
    spline in local coordinates with retraction (Sol\`a); the inner solve is
    unchanged because it never differentiates $\bq$.
  \item \textbf{Contact.} Add $\blambda$ to $w$ with friction-cone/unilateral
    constraints (the inner problem becomes a small SOCP; use an active-set or
    proximal solver; the reduced system is unchanged on the free set).  Treat
    measured GRF as observations of $\blambda$, never as decision variables.
  \item \textbf{Observations.} Markerless 2-D keypoints enter through camera
    projection with robust kernels; marker attachment offsets are geometry
    variables so dimension refinement propagates to the markers.
  \item \textbf{Actuation layer.} Net torques first; an optional declared
    joint-level activation model maps torques to bounded activations with
    first-order dynamics (reported as a model choice, not a measurement).
\end{itemize}

\section{Generalisation to Humanoids and Manipulation}
The kernels are model-agnostic: a humanoid's floating-base rows supply
\eqref{eq:unactuated}, a known payload or the measured ground wrench supplies
the anchor, and the output $(s,\bpi,\bu^\star,\bK)$ is a reference-plus-gains
schedule that whole-body controllers consume directly.  For manipulation the
object plays the club's role (unknown inertia identified from the wrench
transmitted through a known-inertia tool, or vice versa).  Identified
parameters become the nominal model; residual learning or light randomisation
covers what remains, rather than replacing identification.

\section{Review of the DIME Epic and Repository Defects}
\label{sec:dime_review}
The audit of \#11421 (child issues \#11422--\#11437, all merged 2026-10-02..04)
found that the estimator contracts and bookkeeping landed but the estimator
largely did not.  Recorded for correction, with file references at audit
commit \texttt{21a2ed06}:
\begin{itemize}[nosep]
  \item No engine backs any DIME provider; the \texttt{DynamicsProvider}
    protocol lacks derivatives, parameter setters and batching
    (\texttt{dime\_contracts.py:930-956}).
  \item The ``coupled state-control window'' is control-only single shooting
    with hard-coded zero defects and costs (\texttt{dime\_dynamics\_window.py:517-566,604}).
    Partially corrected by \#11554: the least-squares residual now has the fixed
    length $n_u N_{\mathrm{steps}}$ (control effort) $+\ \sum_k \min(\dim\bx^{\mathrm{obs}}_k, n_q)$
    (observations) $+\ (N_{\mathrm{steps}}-1)\,n_u$ (control rate, when weighted)
    $+\ n_u N_{\mathrm{steps}}$ (actuator bounds, when set; one non-negative entry per
    control, zero inside the bounds), independent of the control values and
    enforced by a postcondition (\texttt{expected\_residual\_size}).  A failed
    provider step no longer returns the placeholder residual $10^{8}\mathbf{1}_{100}$
    (whose length differed from the normal path); it raises
    \texttt{DimeResidualEvaluationError}, and the solver reports
    \texttt{failed\_dynamics\_step} with \texttt{success=False}.  The zero
    defects and costs, and the single-shooting structure, remain open
    (\#11545).  Reproduce with
    \texttt{python3 -m pytest tests/unit/estimation/test\_dime\_dynamics\_window.py -k FixedSize}.
  \item MHE never applies its arrival factor (\texttt{moving\_horizon.py:140-213,401,636}).
  \item Learned initialisers, ablations, baseline and replay audit return
    constants (\texttt{dime\_learned\_initializers.py:517-753},
    \texttt{dime\_benchmark\_ablations.py:214-255}, \texttt{dime\_manifest.py:563-629},
    \texttt{dime\_continuous\_replay.py:351-417}); solver-cache timings are fabricated
    and its key omits the targets (\texttt{dime\_solver\_cache.py:293-313,408-418}).
  \item \texttt{multi\_trial.py:379} vs \texttt{:413-426}: locked parameters with a
    prior give a residual/Jacobian row mismatch; its covariance is conditional,
    not marginal.
  \item \texttt{drift\_prediction.py}: covariance propagation omits
    $\partial f/\partial\bx$; a horizon collapses to one step; the default
    selection matrix actuates the root.
  \item \texttt{dime\_manifold.py}: \texttt{local\_coordinates\_jacobian} is the
    identity; the retraction Jacobian is valid only at zero velocity.
  \item ZVCF is defined two incompatible ways (\texttt{ztcf\_zvcf.py:339}
    vs \texttt{counterfactual.py:62,98}); \texttt{\_dynamics\_interface.py:227} has a
    sign error; the GS3DX reference's ZTCF equation carried an undefined
    $\blambda$ under a ``contact free'' heading and integrated the drift along
    the measured state (corrected in that document alongside this one).
    \emph{Resolved (\#11553):} the engine contract documented the drift and
    ZTCF as $\bM^{-1}(\mathbf{C}\bv+\mathbf{g}+\bJ_c^{\!\top}\blambda)$; the bias enters with a
    minus sign, $\bM^{-1}(\bJ_c^{\!\top}\blambda_0-\bh)$, as in
    \eqref{eq:ztcf}.  Engine implementations already computed $-\bM^{-1}\bh$;
    the contract text and its three copies are corrected.
    \texttt{AccelerationDecomposition.zvcf} returned
    $\ba_{\mathrm{grav}}+\bM^{-1}\bB\bu$, the zero-velocity, control-preserved
    acceleration; it now returns \eqref{eq:zvcf} by delegating to
    \texttt{simulation\_backends.ztcf\_zvcf.zvcf\_acceleration}, the single
    implementation, and the old quantity keeps its own name
    (\texttt{zero\_velocity\_control\_preserved}).  On the planar double
    pendulum the two definitions differed by up to
    $73.9\,\mathrm{rad\,s^{-2}}$ at one state; both entry points now agree to
    $10^{-10}$ relative.  Reproduce with
    \texttt{python3 -m pytest tests/unit/motion\_matching/test\_zvcf\_single\_definition.py
    tests/unit/engine\_core/test\_counterfactual\_sign\_convention.py}.
    The duplicate \texttt{DynamicsProvider} protocols remain open.
  \item No inertial regressor existed anywhere in the estimation package.
\end{itemize}

\section{Roadmap}
\label{sec:roadmap}
The follow-up epic (MOSAIC) sequences: engine regressor providers and a
conformance test against this reference; floating-base manifold spline;
contact wrenches with cones and GRF observations; 2-D keypoint observations and
attachment-offset geometry; capture-A/capture-O fits with the gates of
\S\ref{sec:gates}; multi-swing acquisition (the repository holds one trial per
capture); analytic derivatives and MJX batching; activation layer; and
remediation of the defects above.  Promotion criteria are frozen before the
experiments: open-loop replay tolerance, torque RMS against synthetic truth,
and time-to-accepted-match against the IK+smoothing baseline, all measured.

\section{Non-Claims}
This document does not claim real-time performance, a global optimum, muscle
force identification, or validity of any human fit.  Synthetic truth is truth
of the simulator.  Numbers quoted are measured on the stated fixture and
hardware and are reproducible by the stated command.

\begin{thebibliography}{99}
\bibitem{AtkesonAnHollerbach1986} C.~G.~Atkeson, C.~H.~An, J.~M.~Hollerbach. Estimation of inertial parameters of manipulator loads and links. \emph{IJRR} 5(3), 1986.
\bibitem{GautierKhalil1990} M.~Gautier, W.~Khalil. Direct calculation of minimum set of inertial parameters of serial robots. \emph{IEEE T-RA} 6(3), 1990.
\bibitem{Ayusawa2014} K.~Ayusawa, G.~Venture, Y.~Nakamura. Identifiability and identification of inertial parameters using the underactuated base-link dynamics for legged multibody systems. \emph{IJRR} 33(3), 2014.
\bibitem{Wensing2018} P.~M.~Wensing, S.~Kim, J.-J.~E.~Slotine. Linear matrix inequalities for physically consistent inertial parameter identification. \emph{RA-L} 3(1), 2018. arXiv:1701.04395.
\bibitem{Rucker2022} C.~Rucker, P.~M.~Wensing. Smooth parameterization of rigid-body inertia. \emph{RA-L} 7(2), 2022.
\bibitem{Lee2020} T.~Lee, P.~M.~Wensing, F.~C.~Park. Geometric robot dynamic identification: a convex programming approach. \emph{IEEE T-RO} 36(2), 2020.
\bibitem{Traversaro2016} S.~Traversaro, S.~Brossette, A.~Escande, F.~Nori. Identification of fully physical consistent inertial parameters using optimization on manifolds. \emph{IROS} 2016. arXiv:1610.08703.
\bibitem{Wensing2024} P.~M.~Wensing, G.~Niemeyer, J.-J.~E.~Slotine. A geometric characterization of observability in inertial parameter identification. \emph{IJRR}, 2024. arXiv:1711.03896.
\bibitem{GolubPereyra2003} G.~Golub, V.~Pereyra. Separable nonlinear least squares: the variable projection method and its applications. \emph{Inverse Problems} 19, 2003.
\bibitem{Kaufman1975} L.~Kaufman. A variable projection method for solving separable nonlinear least squares problems. \emph{BIT} 15, 1975.
\bibitem{YamashitaFukushima2001} N.~Yamashita, M.~Fukushima. On the rate of convergence of the Levenberg--Marquardt method. \emph{Computing Suppl.} 15, 2001.
\bibitem{More1978} J.~J.~Mor\'e. The Levenberg--Marquardt algorithm: implementation and theory. \emph{Numerical Analysis}, LNM 630, 1978.
\bibitem{Nielsen1999} H.~B.~Nielsen. Damping parameter in Marquardt's method. IMM-REP-1999-05, DTU, 1999.
\bibitem{RemyThelen2009} C.~D.~Remy, D.~G.~Thelen. Optimal estimation of dynamically consistent kinematics and kinetics for forward dynamic simulation of gait. \emph{J.~Biomech.\ Eng.} 131(3), 2009.
\bibitem{Falisse2019} A.~Falisse \emph{et al.} Rapid predictive simulations with complex musculoskeletal models suggest that diverse healthy and pathological human gaits can emerge from similar control strategies. \emph{J.~R.~Soc.\ Interface} 16, 2019.
\bibitem{Werling2023} K.~Werling \emph{et al.} AddBiomechanics: automating model scaling, inverse kinematics, and inverse dynamics from human motion data through sequential optimization. \emph{PLOS ONE}, 2023.
\bibitem{Dembia2020} C.~L.~Dembia \emph{et al.} OpenSim Moco: musculoskeletal optimal control. \emph{PLOS Comput.\ Biol.}, 2020.
\bibitem{Giftthaler2018} M.~Giftthaler \emph{et al.} A family of iterative Gauss--Newton shooting methods for nonlinear optimal control. \emph{IROS} 2018. arXiv:1711.11006.
\bibitem{Jallet2025} W.~Jallet \emph{et al.} PROXDDP: proximal constrained trajectory optimization. \emph{IEEE T-RO}, 2025.
\bibitem{Singh2022} S.~Singh, R.~P.~Russell, P.~M.~Wensing. Efficient analytical derivatives of rigid-body dynamics using spatial vector algebra. \emph{RA-L} 7(2), 2022.
\bibitem{Sarkka2023} S.~S\"arkk\"a, \'A.~F.~Garc\'ia-Fern\'andez. Temporal parallelization of dynamic programming and linear quadratic control. \emph{IEEE TAC}, 2023.
\bibitem{Sola2018} J.~Sol\`a, J.~Deray, D.~Atchuthan. A micro Lie theory for state estimation in robotics. arXiv:1812.01537, 2018.
\bibitem{Bristow2006} D.~A.~Bristow, M.~Tharayil, A.~G.~Alleyne. A survey of iterative learning control. \emph{IEEE Control Systems Magazine} 26(3), 2006.
\bibitem{deLeva1996} P.~de~Leva. Adjustments to Zatsiorsky--Seluyanov's segment inertia parameters. \emph{J.~Biomech.} 29(9), 1996.
\bibitem{AffineDriftZTCF} AffineDrift. The zero-torque counterfactual. \url{https://www.affinedrift.com/articles/zero-torque-counterfactual.html}.
\bibitem{DIME} UpstreamDrift. Epic \#11421, Dynamics-Informed Mocap Matching With ZTCF Prediction and Continuous Forward Replay, and children \#11422--\#11437.
\end{thebibliography}
\end{document}
